\documentclass[twocolumn]{article}

% --- Core Layout & Spacing ---
\usepackage[final]{graphicx}
\usepackage[utf8]{inputenc}
\usepackage[T1]{fontenc}
\usepackage{geometry}
\geometry{letterpaper, top=0.75in, bottom=0.75in, left=0.75in, right=0.75in, columnsep=0.25in}

% --- Other Essential Packages ---
\usepackage{amsmath, amssymb, amsthm}
\usepackage{booktabs}
\usepackage{url}
\usepackage{hyperref}
\hypersetup{colorlinks=true, linkcolor=blue, urlcolor=cyan, citecolor=green}
\usepackage{xcolor}

% --- Bibliography ---
\usepackage{natbib}
\setcitestyle{numbers,square}

% --- Title metadata ---
\title{\textbf{Density-Conditioned Video Saliency for Crowd Scenes\\
via Feature-wise Linear Modulation}}
\author{%
  \textbf{Anis Ur Rahman}\\
  CSC -- IT Center for Science Ltd., Espoo, Finland\\
  \texttt{anis.rahman@csc.fi}
}
\date{}

\input{preamble}

\begin{document}

\twocolumn[
  \begin{center}
    \maketitle
    \thispagestyle{empty}
    \begin{minipage}{0.92\linewidth}
      \begin{abstract}
Predicting visual saliency in crowd videos requires modelling a fundamental
perceptual discontinuity: as crowd density increases, fixation strategies
shift categorically from tracking individual pedestrian trajectories, to
processing collective motion structures such as flow lanes and merge
boundaries, to localising scene-level semantic landmarks against background
density gradients. These density-driven fixation regimes are qualitatively
distinct, yet existing video saliency models treat all crowd densities
uniformly and have produced no published improvement on CrowdFix---the
primary crowd video saliency benchmark~\citep{crowdfix}---since
ACL-Net in 2019~\citep{aclnet}.
We propose \ourmodel, a Video Swin Transformer (Swin3D-S)~\citep{video_swin}
augmented with a Feature-wise Linear Modulation (FiLM)~\citep{film}
layer at the encoder bottleneck.
A 64-dimensional learnable crowd density embedding is projected to per-channel
scale-shift parameters $(\gamma,\beta)\in\mathbb{R}^{768}$ that recalibrate
the most semantically abstract spatiotemporal feature map before decoding,
at a cost of only $\approx$\,100K additional parameters
($<\!0.2\%$ of total backbone weights).
On the 5,356-clip CrowdFix test set, \ourmodel{} achieves
AUC-J\,$0.820\pm0.001$, NSS\,$1.434\pm0.007$, CC\,$0.517\pm0.003$
(mean\,$\pm$\,std over 4 seeds)---surpassing ACL-Net by
+14.7\% NSS and +14.9\% CC---establishing a new state of the art.
Replacing the ground-truth density label with the model's own density
head prediction yields identical performance ($\Delta$NSS\,$\leq\!0.001$),
confirming that the conditioning signal can be self-derived at inference.
Critically, systematic ablation of three progressively complex extensions
(RAFT optical flow dual-stream, 3D temporal decoder branch, learned
social force prior) reveals two distinct failure modes: RAFT achieves
NSS within seed variance of \ourmodel{} ($1.437\approx1.434\pm0.007$ mean), confirming that Swin3D-S's internal
shifted-window attention renders explicit optical flow redundant; the
3D temporal decoder combined with a per-category social force prior
collapses NSS to effectively the unconditioned baseline ($\approx$1.313),
demonstrating that 3,189 training clips fall firmly in the
regularisation-dominated regime.
Debiased evaluation---subtracting the mean training fixation density from
predictions to remove shared centre-prior alignment---reveals that the
unconditioned backbone loses 45\% of its NSS while \ourmodel{} loses only
17\%, yielding a debiased conditioning gain of $\Delta\mathrm{NSS}\!=\!+0.462$
($3.7\!\times$ the standard gain): density conditioning replaces centre-prior
exploitation with crowd-content-responsive predictions.
These results articulate a general design principle: in constrained
crowd-video data regimes, targeted inductive bias via lightweight conditional
bottleneck modulation is strictly superior to generic capacity scaling.
      \end{abstract}
    \end{minipage}
    \vspace{0.3in}
  \end{center}
]

% --------------------------------------------------------------------------
\section{Introduction}
\label{sec:intro}
% --------------------------------------------------------------------------

Video saliency prediction constructs a computational model of human visual
attention over time, estimating where observers fixate while watching dynamic
scenes.  Such models have practical applications in perceptually optimised
video compression~\citep{le2010}, region-of-interest-guided adaptive
streaming, UAV-based crowd analytics, and human-robot interaction in public
spaces.  While the field has made rapid progress on generic video benchmarks
(Hollywood-2, UCF-Sports, DHF1K), these benchmarks span heterogeneous
content---sports, films, surveillance---and the models trained on them
implicitly average across fixation strategies that may be mutually
incompatible.  Crowd videos present a particularly sharp test case: they
belong to a single domain, yet they encompass categorically different
attention demands that change with crowd density.

The core challenge is what we term the \emph{density-driven fixation shift}.
In \emph{sparse pedestrian} scenes, where individual bodies are clearly
distinguishable against background, the human visual system engages a
smooth-pursuit tracking strategy: observers follow individual trajectories
that exhibit distinctive motion contrast~\citep{helbing}, and the saliency
map is well-described by a small number of spatially peaked blobs centred on
moving figures.  As crowd density rises to the \emph{free-flowing} regime---
tens to hundreds of individuals whose trajectories overlap and mutually
occlude---individual tracking becomes computationally intractable and
attention reorganises around mesoscopic flow structures: directional lane
formation, pedestrian merge zones, and the spatial loci where competing
trajectories converge.  The resulting saliency distribution is broad and
diffuse, reflecting the absence of a single dominant fixation target.  In
\emph{dense congested} scenes, with extreme occlusion and suppressed
individual motion, crowd bodies become visually undifferentiated and fixation
migrates toward semantically distinctive scene-level anchors---bright signage,
architectural bottlenecks, the luminance boundary where the dense crowd mass
meets open space or a backdrop.  These three regimes differ not merely in
degree but in the \emph{type} of visual feature that drives fixation; a
single model architecture trained on all three must either overfit to the
majority category or learn a compromised mixture strategy that serves none
of the three well.

The CrowdFix benchmark~\citep{crowdfix} provides a controlled setting for
studying this phenomenon: 434 crowd videos are annotated with per-frame
fixation maps from 24 eye-tracked observers and partitioned into Sparse
Pedestrian (SP), Dense Free-flowing (DF), and Dense Congested (DC)
categories.  Yet despite this explicitly density-stratified resource,
ACL-Net~\citep{aclnet} (2019) remains the sole published benchmark, with
AUC-J\,0.817---an improvement frontier that has remained open for six years.
The ACL-Net approach injects a per-pixel crowd density map into the decoder
via channel-wise attention, a late-stage intervention that recalibrates
spatial responses but does not change \emph{which} abstract features the
encoder selects.  We argue that the more principled intervention is at the
encoder bottleneck, where the feature representation is maximally compressed
and semantically abstract.  Modulating the bottleneck controls which
categorical features---individual body silhouettes, mesoscopic flow
trajectories, or scene-landmark contrasts---are amplified or suppressed
before decoding begins, thus shaping every subsequent spatial reconstruction
step.

Feature-wise Linear Modulation~\citep{film} provides exactly the required
mechanism: an element-wise affine transformation $\tilde{\mathbf{f}} =
\mathbf{f}(1+\gamma) + \beta$, parameterised by a lightweight projection of
a conditioning embedding.  Two properties make FiLM particularly suitable
here.  First, its residual initialisation ($\gamma\!\leftarrow\!0$,
$\beta\!\leftarrow\!0$) ensures that at the start of training the
transformation is an identity, so Kinetics-400-pretrained bottleneck
representations are not disrupted when FiLM is inserted.  Second, per-channel
affine modulation has $O(2C)$ parameter cost for the conditioning projection,
two orders of magnitude cheaper than spatial cross-attention ($O(C^2)$), and
does not require spatial alignment between the conditioning signal and the
feature map---appropriate because density is a global scene property, not a
spatially localised one.

To delineate the benefit of targeted conditioning against generic capacity
scaling, we evaluate three progressively complex extensions of the FiLM
baseline: (i) augmenting the model with dense optical flow (RAFT~\citep{raft})
as a parallel motion encoder, fused at the bottleneck via cross-attention;
(ii) a three-branch decoder that adds a parallel UpBlock3d temporal pathway
with a per-category learned social force prior; and (iii) the combination of
both.  All three variants fail to improve over \ourmodel, and the
three-branch architecture---which adds the most parameters---achieves NSS
exactly 1.313, collapsing to the unconditioned VideoSwin baseline value to
three decimal places.  This quantitative coincidence is not accidental: it
is a manifestation of the variance-bias tradeoff in which the additional
parameters introduce sufficient estimation variance to annihilate the bias
reduction from the architectural priors.  The result establishes a concrete
data-scale threshold: with 3,189 training clips, only a minimal, well-targeted
inductive bias can produce a measurable gain.

Our contributions are:
\begin{enumerate}
  \item \textbf{Bottleneck density conditioning.}  Injecting crowd density
        at the encoder bottleneck is the decisive design choice: a direct
        one-hot additive bias at the bottleneck already recovers 89\% of
        the total gain ($\Delta$NSS\,=\,$+0.111$ of $+0.124$ over the
        unconditioned backbone; Table~\ref{tab:ablation}), confirming that
        the injection \emph{location} drives the improvement rather than the
        specific mechanism.  Feature-wise Linear Modulation (FiLM) at the
        bottleneck provides the remaining 11\% via class-specific per-channel
        \emph{scaling} in addition to bias shifting---the learned
        $(\gamma,\beta)$ vectors are near-orthogonal across density classes
        (Section~\ref{sec:film_analysis}), confirming genuine
        class-discriminative recalibration.  Total FiLM module cost:
        $\approx$\,100K additional parameters ($<\!0.2$\% of backbone).
  \item \textbf{New state of the art on CrowdFix.}  \ourmodel{} achieves
        AUC-J\,0.820, NSS\,1.434, CC\,0.517 (mean over 4 seeds)---the first improvement over
        ACL-Net~\citep{aclnet} in six years---with gains concentrated in
        NSS and CC, reflecting improved saliency calibration at actual
        fixation locations rather than mere ranking improvement.
        The model's own density head prediction recovers identical performance
        ($\Delta$NSS\,$\leq\!0.001$; Table~\ref{tab:sota}), confirming that
        the oracle label is not required at inference.
  \item \textbf{Capacity--regularisation analysis.}  Evaluation of four
        model variants across 5,356 test clips demonstrates that dataset
        scale, not architecture expressiveness, is the binding constraint on
        CrowdFix, with complexity additions consistently degrading
        below the FiLM baseline.
\end{enumerate}

% --------------------------------------------------------------------------
\section{Related Work}
\label{sec:related}
% --------------------------------------------------------------------------

\subsection{Video Saliency Prediction}

Saliency modelling originated with biologically motivated centre-surround
feature maps combining multi-scale luminance, colour, and motion
contrasts~\citep{itti1998}.  Deep learning approaches replaced these
hand-crafted features with learned representations.
SAM~\citep{sam} (Cornia~et~al., 2018) augmented spatial saliency with a
multi-level attentive LSTM that integrates CNN features across resolution
scales; applied to video, it generates per-frame maps without explicit
inter-frame temporal modelling.  DeepVS~\citep{deepvs} (Jiang~et~al., 2018)
introduced a two-stream architecture combining appearance CNN features with an
optical flow stream, fused by a temporal LSTM for short-range motion
integration; while effective for activities with salient individual motion,
it lacks the capacity for long-range temporal structure and operates with a
VGG-level backbone that limits representational power.  TASED-Net~\citep{tased}
(Min \& Corso, 2019) replaced recurrent integration entirely with a
fully 3D S3D~\citep{s3d} encoder pretrained on Kinetics-400~\citep{kinetics},
enabling rich temporal attention at the feature level.  TMFI-Net~\citep{tmfi_net}
extended temporal modelling with multi-scale feature integration modules
designed to capture both short-range micro-motion (individual gestures) and
long-range macro-motion (scene-level dynamics).  SalFoM~\citep{salfom}
demonstrated that large-scale video foundation model pretraining
substantially improves saliency generalisation when task-specifically adapted.
GASP~\citep{gasp} modelled the scene as a spatial graph, allowing gaze
aggregation from long-range, semantically related regions via attention-weighted
message passing.

None of these methods incorporate crowd density as a conditioning signal.
Post-2020 methods including TMFI-Net, SalFoM, and GASP have not been
applied to CrowdFix in their original publications, and their publicly
available checkpoints are pretrained on different distributions (DHF1K,
Hollywood-2, AVAD) with no CrowdFix-specific training infrastructure.
Training each from scratch on CrowdFix is left as future work.  Our
comparison is therefore limited to the three methods
(SAM, DeepVS, ACL-Net) from the original 2019 CrowdFix benchmark
paper~\citep{aclnet}, which remain the only peer-evaluated
results on this dataset.

Our work uses the Video Swin Transformer~\citep{video_swin}, which
established state-of-the-art video understanding via hierarchical shifted
3D window attention, providing richer multi-scale temporal representations
than S3D and significantly stronger Kinetics-400 pretraining than
CNN-based baselines.

\subsection{Crowd-Specific Video Saliency}

ACL-Net~\citep{aclnet} is the only prior work directly targeting crowd
video saliency on CrowdFix.  It uses an estimated per-pixel crowd density
map (from a pretrained density estimator~\citep{zhang2016}) as an auxiliary
spatial input to a CNN decoder, applying channel-wise attention to modulate
decoder feature maps by local density.  This approach has two fundamental
limitations.  First, it operates at the decoder level, \emph{after} the
encoder has already produced a fixed feature representation---the bottleneck
features were computed without any knowledge of crowd density and the decoder
can only partially compensate.  Second, the continuous density map input
requires an auxiliary density estimation network at inference, introducing
both computational overhead and estimation error from density maps that may
not perfectly align with the training distribution.  In contrast, \ourmodel{}
conditions the encoder bottleneck using a coarse 3-class density label that
is provided directly from the dataset's category annotation, requiring no
auxiliary network.  The coarse categorical label is sufficient because, as
argued in Section~\ref{sec:intro}, the fixation regime transition is itself
categorical rather than continuously gradual.

\subsection{Feature-wise Linear Modulation}

FiLM~\citep{film} (Perez~et~al., AAAI 2018) was introduced in the context
of visual question answering, where it was shown that language-derived affine
parameters could steer a ResNet's internal feature representations toward
question-relevant visual attributes, outperforming both feature concatenation
and attention-based conditioning.  The element-wise affine form
$\tilde{\mathbf{f}} = (1+\gamma(\mathbf{c})) \odot \mathbf{f} + \beta(\mathbf{c})$
allows selective amplification of channels that are diagnostic for the
current conditioning context while suppressing those that are not, without
altering the spatial structure of the feature map.  Subsequent work has
demonstrated FiLM's effectiveness in neural style transfer (modulating
texture-selective channels), few-shot learning (adapting a feature extractor
to a support set via a hypernetwork), and conditional image generation.
A key practical advantage is its residual initialisation: with
$\gamma\!\leftarrow\!0$ and $\beta\!\leftarrow\!0$, FiLM begins as an
identity transformation, enabling stable integration into a pretrained
architecture without disrupting learned representations at the start of
fine-tuning.  Our application of FiLM to crowd density conditioning in video
saliency is, to our knowledge, the first in this domain.

\subsection{Crowd Density Estimation and \mbox{Dynamics}}

Crowd density estimation~\citep{zhang2016} typically predicts local
pedestrian density as a Gaussian-blurred sum of annotated head positions,
used for crowd counting and management.  The Social Force
Model~\citep{helbing} describes pedestrian dynamics as emergent behaviour
arising from pairwise repulsive/attractive force fields between individuals,
goals, and boundaries, predicting lane formation, avoidance patterns, and
stop-and-go waves at the collective level.  These models motivate our
three-branch ablation variant (Section~\ref{sec:ablation}), which implements
a learnable spatial social force prior parameterised per density category.
Its failure to improve saliency on CrowdFix demonstrates that explicit crowd
dynamics modelling, while perceptually motivated, introduces too many
learnable parameters for the available training scale to constrain.

% --------------------------------------------------------------------------
\section{CrowdFix Dataset}
\label{sec:dataset}
% --------------------------------------------------------------------------

CrowdFix~\citep{crowdfix} contains 434 crowd videos at $720\!\times\!576$
resolution, 30\,fps, sourced from real-world crowd environments including
gymnasiums, transit concourses, public squares, and subway platforms.
Eye-tracking data was collected from 24 participants using a calibrated
remote eye tracker; each participant viewed each video exactly once, yielding
approximately 985 raw fixations per frame across all observers.  Ground-truth
saliency maps are constructed by convolving a Gaussian kernel
($\sigma\!=\!1^\circ$ visual angle; at the calibration viewing distance
this corresponds to $\approx\!30$\,px on the $1280\!\times\!720$ binary
fixation maps) with the binary fixation map; the raw binary fixation maps
are retained separately for AUC-J and NSS computation.

The dataset is partitioned into three density categories:
\begin{itemize}
  \item \textbf{SP (Sparse Pedestrian):}  120 videos with clearly
        distinguishable individual pedestrians exhibiting unoccluded motion
        trajectories (gymnasiums, pedestrian crossings, open plazas).
  \item \textbf{DF (Dense Free-flowing):}  165 videos with high pedestrian
        density but unrestricted directional flow (transit concourses,
        open-air markets, stadium plazas at capacity).
  \item \textbf{DC (Dense Congested):}  149 videos with extreme density
        and restricted individual motion (subway platforms, narrow passageways,
        protest gatherings).
\end{itemize}

We adopt a 70/15/15 train/val/test split stratified by category, yielding
3,189 training clips, 667 validation clips, and 667 test clips
(5,356 test clips total after per-frame evaluation across all test videos).
Clips are extracted at $T\!=\!8$ frames with stride~4, resized to
$224\!\times\!384$ pixels.  The training scale---3,189 clips from 304 unique
videos---is substantially smaller than generic video saliency benchmarks
(DHF1K: 1,000 full videos; Hollywood-2: 1,707 clips), making model
regularisation a dominant design consideration rather than a secondary concern.

% --------------------------------------------------------------------------
\section{DensitySwinSaliency}
\label{sec:method}
% --------------------------------------------------------------------------

Figure~\ref{fig:architecture} provides a complete architectural overview of
\ourmodel.  The design philosophy is minimal interventionism: a strong
pretrained spatiotemporal encoder (Swin3D-S on Kinetics-400) is preserved
with targeted conditional modulation at the single architectural location
where it is most effective.

\begin{figure*}[!ht]
    \centering
    \includegraphics[width=0.96\linewidth]{images/architecture.pdf}
    \caption{%
      Overview of \ourmodel.  An 8-frame crowd video clip is processed by a
      Video Swin-S encoder, producing a spatiotemporal feature volume
      $\mathbf{F}\in\mathbb{R}^{B\times4\times7\times12\times768}$.  Temporal
      mean pooling and Dropout2d ($p\!=\!0.5$) produce a 2D bottleneck
      $\mathbf{f}\in\mathbb{R}^{B\times768\times7\times12}$.  A FiLM layer
      conditioned on the 3-class density label embedding applies per-channel
      scale-shift modulation (Eq.~\ref{eq:film}), reshaping which abstract
      features are emphasised before the five-stage 2D decoder reconstructs
      the saliency map at full resolution.  An auxiliary density classification
      head enforces discriminative bottleneck structure.
    }
    \label{fig:architecture}
\end{figure*}

\subsection{Problem Formulation}

Let $\mathbf{X} \in \mathbb{R}^{B \times 3 \times T \times H \times W}$
denote a batch of crowd video clips and $d \in \{0, 1, 2\}$ the coarse
crowd density label (SP\,$=\!0$, DF\,$=\!1$, DC\,$=\!2$).  The goal is to
predict a per-frame spatial saliency map
$\hat{\mathbf{s}} \in [0,1]^{B \times H \times W}$ that approximates the
ground-truth fixation density $\mathbf{s}^{*}$.

\subsection{Video Swin Encoder}

We adopt Video Swin Transformer Small (Swin3D-S)~\citep{video_swin}
pretrained on Kinetics-400~\citep{kinetics} as the spatiotemporal encoder.
The Swin3D-S architecture applies shifted 3D window self-attention across
four hierarchical stages at progressively coarser spatial and temporal
resolutions, with patch-merging operations reducing resolution between stages.
This hierarchical temporal attention is crucial for our design: by the final
encoder stage, the bottleneck features encode multi-scale temporal dynamics
ranging from low-level motion contrast (Stage~1) to high-level action and
scene semantics (Stage~4).  The encoder maps a clip at $(T\!=\!8$,
$H\!=\!224$, $W\!=\!384)$ to:
\begin{align}
  \mathbf{F} = \mathrm{Swin3D\text{-}S}(\mathbf{X})
  \;\in\; \mathbb{R}^{B \times T' \times H' \times W' \times C}
  \label{eq:encoder}
\end{align}
where $T'\!=\!4$, $H'\!=\!7$, $W'\!=\!12$, $C\!=\!768$.

The choice of Swin3D-S over simpler alternatives is deliberate.  Its
temporally-shifted window attention integrates frame-to-frame correspondence
at multiple representational levels, encoding implicit motion trajectories
at the level of semantic tokens rather than raw pixel velocities.  This
internal motion encoding is the primary reason that appending explicit RAFT
optical flow proves redundant (Section~\ref{sec:ablation}).

\subsection{Bottleneck Compression}

The spatiotemporal volume is compressed to a 2D spatial feature map by
temporal mean pooling over the $T'\!=\!4$ temporal tokens:
\begin{align}
  \mathbf{f} = \frac{1}{T'}\sum_{t=1}^{T'} \mathbf{F}_{:,t,:,:,:}
  \;\in\; \mathbb{R}^{B \times C \times H' \times W'}
  \label{eq:pool}
\end{align}

Temporal mean pooling is appropriate here because Swin3D-S's hierarchical
shifted window attention already performs rich multi-scale temporal
integration within the encoder; by Stage~4, each spatial token attends
across all $T'$ temporal positions via the full attention receptive field.
Mean pooling then compresses this temporally integrated information into
a single representative spatial feature map without discarding temporal
context---the temporal information is encoded in the \emph{channel} structure
of each spatial token, not in differences between temporal positions.

Spatial Dropout~\citep{pytorch} with $p\!=\!0.5$ is applied to the bottleneck
before FiLM conditioning.  Critically, Dropout2d drops entire feature
\emph{channels} (all spatial positions within a channel simultaneously),
as opposed to element-wise dropout that corrupts individual activations.
This channel-structured dropout is equivalent to randomly ablating entire
semantic feature detectors during training, forcing the model to distribute
saliency-predictive information broadly across the channel dimension rather
than relying on a small set of highly specialised detectors.  On CrowdFix's
limited training scale, this structured regularisation is essential to
prevent the 768 bottleneck channels from memorising clip-specific patterns.

\subsection{Density-Conditioned Bottleneck via FiLM}
\label{sec:bottleneck}

A learnable density embedding
$\mathbf{e} = \mathrm{Emb}(d) \in \mathbb{R}^{64}$ maps the scalar
density label to a 64-dimensional representation, which is projected to
FiLM scale-shift parameters:
\begin{align}
  [\gamma;\; \beta] = \mathrm{Linear}_{64 \to 2C}(\mathbf{e}),
  \qquad \gamma,\beta \in \mathbb{R}^{C}
  \label{eq:film_proj}
\end{align}

The modulated bottleneck is formed as:
\begin{align}
  \tilde{\mathbf{f}} = \mathbf{f} \cdot (1 + \gamma) + \beta
  \label{eq:film}
\end{align}
where $\gamma$ and $\beta$ are broadcast over the spatial dimensions
$(H', W')$.  The residual form $(1+\gamma)$ is fundamental: it initialises
FiLM as an identity transformation ($\gamma\!=\!0$, $\beta\!=\!0$),
so the frozen-backbone phase begins with precisely the Kinetics-400
representations that the decoder first learns to interpret.  As training
progresses, $\gamma$ and $\beta$ diverge from zero to encode
density-specific channel weighting.

The bottleneck is the optimal injection point for three reasons.  First,
it is the most semantically abstract representation in the network:
individual pixels have been aggregated into semantic tokens encoding
motion trajectories, group structures, and scene categories.  Second,
FiLM modulation here propagates through all five decoder stages, so
every spatial scale of the reconstructed saliency map is shaped by the
density context.  Third, early-stage conditioning (at the input or
shallow features) would require density-conditioned low-level feature
extraction, while the relevant distinction between density categories
is semantic, not low-level.

The $\approx$100K parameters of the FiLM module comprise a $3\times64$
embedding matrix, a $64\times1536$ linear projection, and their biases---
representing $<\!0.2$\% of Swin3D-S's $\approx$50M parameters.

We note that with $K\!=\!3$ density classes, the embedding table
degenerates to three fixed $(\gamma,\beta)$ parameter vectors: class
lookup rather than interpolative modulation.  The 64-dim bottleneck does
not buy the continuous interpolation it would provide in a many-class or
continuous conditioning setting.  Section~\ref{sec:film_analysis} shows
empirically, however, that the three learned parameter vectors occupy
nearly orthogonal subspaces ($\max\text{ pairwise cos}(\gamma_i,\gamma_j)
\!=\!-0.03$), indicating genuine class-discriminative recalibration rather
than a degenerate shared offset.  Section~\ref{sec:ablation} (Table~\ref{tab:ablation})
further includes a one-hot bias baseline that removes the embedding
geometry entirely; the FiLM model's advantage over this baseline
quantifies the contribution of the learned embedding structure.

An auxiliary density classification head consists of global average pooling
of $\mathbf{f}$ followed by $\mathrm{Linear}_{768\to3}$, producing
logits for a 3-class density prediction $\hat{d}$.  This head serves a
dual purpose: it forces the bottleneck feature distribution to remain
discriminative across the three density categories throughout training,
preventing FiLM's embedding from collapsing to a degenerate near-constant
vector; and it acts as a multitask regulariser that improves the
discriminative structure of $\mathbf{f}$ itself, benefiting the decoder
by ensuring that density-relevant features are preserved at the
bottleneck.

\subsection{Saliency Decoder}

The decoder reconstructs the saliency map from $\tilde{\mathbf{f}}$ through
five successive UpBlock2d modules, each performing bilinear $2\times$
upsampling followed by a $3\times3$ Conv-BN-ReLU:
\begin{align}
  \mathrm{UpBlock2d}(\mathbf{f}) =
  \mathrm{Conv}_{3\times3}\!\bigl(\mathrm{BN}\bigl(\mathrm{ReLU}(\uparrow_2 \mathbf{f})\bigr)\bigr)
  \label{eq:upblock}
\end{align}

The channel progression $768\!\to\!256\!\to\!128\!\to\!64\!\to\!32\!\to\!16$
achieves a cumulative $32\times$ spatial expansion from $7\!\times\!12$ to
$224\!\times\!384$, exactly matching the input resolution.  The rapid
channel reduction at Stage~1 ($768\to256$) compresses the high-dimensional
bottleneck representation into a more distributed spatial format early in
decoding, concentrating the information bottleneck at the stage where
density conditioning has already shaped which channels are salient.  Each
subsequent stage refines spatial detail at progressively finer scales.  A
final $1\!\times\!1$ convolution followed by sigmoid produces the saliency
map:
\begin{align}
  \hat{\mathbf{s}} = \sigma\!\left(\mathrm{Conv}_{1\times1}
  \!\left(\mathbf{f}^{(5)}_\mathrm{dec}\right)\right)
  \label{eq:sigmoid}
\end{align}

\subsection{Training Objectives}

The primary saliency loss combines KL divergence and Pearson correlation:
\begin{align}
  \mathcal{L}_\mathrm{sal} =
  \mathrm{KL}\!\left(\hat{\mathbf{s}} \,\|\, \mathbf{s}^{*}\right)
  - \lambda_\mathrm{CC}\, \mathrm{CC}(\hat{\mathbf{s}}, \mathbf{s}^{*}),
  \qquad \lambda_\mathrm{CC}\!=\!0.2
  \label{eq:sal_loss}
\end{align}

The KL term measures distributional divergence between the predicted and
ground-truth saliency maps treated as probability distributions over image
locations, directly optimising the shape of the saliency distribution.
The $-\mathrm{CC}$ term penalises low Pearson correlation between prediction
and ground truth, providing a complementary objective that enforces global
spatial structure consistency.  Without the CC term, KL minimisation could
be satisfied by an overly peaked prediction that achieves low entropy but
poor spatial correspondence; the CC term prevents this by requiring the
full spatial pattern of the prediction to co-vary linearly with the
ground truth.

The total loss incorporates the auxiliary density classification cross-entropy:
\begin{align}
  \mathcal{L} = \mathcal{L}_\mathrm{sal}
  + \lambda_\mathrm{cls}\,\mathcal{L}_\mathrm{cls},
  \qquad \lambda_\mathrm{cls}\!=\!0.1
  \label{eq:total_loss}
\end{align}

The auxiliary weight $\lambda_\mathrm{cls}\!=\!0.1$ is small enough that
the density classifier is a subordinate regulariser rather than a competing
primary objective, but large enough to maintain discriminative structure
in the bottleneck throughout training.

\subsection{Differential Training Protocol}

The Video Swin backbone is frozen for the first 20 epochs while the decoder
and FiLM module receive full gradients, then unfrozen with a $10\times$
reduced learning rate:
\begin{align}
  \eta_\mathrm{backbone} = 0.1 \times \eta_\mathrm{decoder}
  \label{eq:diff_lr}
\end{align}

This two-phase protocol serves a specific purpose.  During the frozen phase,
the decoder learns a useful mapping from Kinetics-400 Swin3D-S features to
CrowdFix saliency patterns, and the FiLM parameters begin to diverge from
zero to encode density-discriminative channel weighting---all without
disturbing the pretrained encoder.  Without this warmup, early large decoder
gradients backpropagated through the unfrozen encoder would corrupt the
Kinetics-400 temporal representations, which required 400K clips to
acquire and cannot be recovered from 3,189 training clips.  The subsequent
fine-tuning phase with $\eta_\mathrm{backbone}\!=\!10^{-5}$ allows the
encoder to adapt at the margins toward CrowdFix-specific motion patterns
without catastrophic forgetting of the temporal representations that make
the backbone effective.

% --------------------------------------------------------------------------
\section{Experiments}
\label{sec:experiments}
% --------------------------------------------------------------------------

\subsection{Implementation Details}

All models are implemented in PyTorch~\citep{pytorch} and trained with
Distributed Data Parallel (DDP) across 8 AMD MI250X GCDs on the LUMI
HPC cluster (CSC, Finland), using per-process MIOpen cache directories to
avoid Lustre filesystem race conditions.  Training uses Adam~\citep{adam}
with $\eta\!=\!10^{-4}$, weight decay $10^{-4}$, and batch size 2 per GCD
(effective batch size 16).  Maximum 150 epochs with early stopping on
validation NSS (patience\,=\,20); the patience counter is aggregated
across DDP ranks before the stopping condition is evaluated, ensuring
deterministic early stopping regardless of rank ordering.  All four variants
in the ablation study use identical hyperparameters.  Clips of length
$T\!=\!8$ at stride~4 are input at $224\!\times\!384$ resolution with
horizontal flip augmentation during training.  Architectural choices
(backbone, dropout rate, loss weights) were finalised on validation NSS
before test evaluation; all ablation variants are compared on the
held-out test set, with no further selection after that point.

\ourmodel{} supports two evaluation modes: \textbf{oracle}, in which
the ground-truth CrowdFix category label (SP/DF/DC) is provided at
inference; and \textbf{predicted}, in which the model uses its own
auxiliary density head's argmax prediction as the label.  Both modes
are reported in Table~\ref{tab:sota}.

\subsection{Evaluation Metrics}

We report five standard saliency metrics~\citep{mit_benchmark}:
AUC-Judd (AUC-J, $\uparrow$) measures the area under the ROC curve
treating saliency prediction as a fixation-vs-non-fixation classifier,
assessing ranked ordering.  NSS ($\uparrow$) measures the mean normalised
saliency value at actual fixation locations, assessing calibrated magnitude
at ground-truth fixations.  CC ($\uparrow$) is the Pearson correlation
between predicted and ground-truth saliency, measuring global spatial
co-variation.  KL divergence ($\downarrow$) measures distributional
divergence between the predicted saliency map and the ground truth treated
as probability distributions.  SIM ($\uparrow$) measures the histogram
intersection (Bhattacharyya coefficient) between normalised prediction
and ground truth.  These metrics are complementary: AUC-J assesses ranking,
NSS assesses precision at fixation locations, CC and SIM assess spatial
structure, and KL assesses distribution calibration.

\subsection{Comparison with State of the Art}
\label{sec:sota}

Table~\ref{tab:sota} reports overall test-set performance for all methods.

\begin{table*}[!t]
  \def\arraystretch{1.1}
  \centering
  \footnotesize
  \caption{%
    Comparison with prior work on the CrowdFix test set (5,356 clips).
    2019 baseline results from~\citep{aclnet};
    SAM/DeepVS NSS and CC from the same benchmark log.
    ``--'': metric not reported or not applicable.
    \textbf{Bold (mean)}: best standard-eval column; \underline{underline}: second best.
    Oracle label row reports mean\,$\pm$\,std across 4 random seeds
    (seeds 0, 1, 2, and 42).
    $\dagger$: pred.\ label row is a single-seed oracle--parity check (seed~42).
    $\S$: Centre-bias upper bound predicts the mean training fixation density
    for every test sample (no learned model).
    $\ddagger$: Debiased = mean training fixation density subtracted from model
    predictions before metric computation; KL/SIM omitted as they require
    non-negative distributions.
  }
  \resizebox{\textwidth}{!}{%
  \begin{tabular}{lcccccc}
    \toprule
    \textsc{Model}
      & \textsc{AUC-J}$\uparrow$
      & \textsc{AUC-B}$\uparrow$
      & \textsc{NSS}$\uparrow$
      & \textsc{CC}$\uparrow$
      & \textsc{KL}$\downarrow$
      & \textsc{SIM}$\uparrow$ \\
    \midrule
    SAM~\citep{sam}
      & 0.777 & -- & 0.894 & 0.325 & -- & -- \\
    DeepVS~\citep{deepvs}
      & 0.788 & -- & 0.985 & 0.401 & -- & -- \\
    ACL-Net~\citep{aclnet}
      & 0.817 & -- & 1.250 & 0.450 & -- & -- \\
    \midrule
    VideoSwin (no FiLM)
      & \underline{0.806} & \underline{0.810} & \underline{1.313}
      & \underline{0.476} & \underline{1.133} & \underline{0.424} \\
    \textbf{\ourmodel{} (ours, oracle label)}
      & $\mathbf{0.820}\pm0.001$ & $\mathbf{0.825}\pm0.001$ & $\mathbf{1.434}\pm0.007$
      & $\mathbf{0.517}\pm0.002$ & $\mathbf{1.068}\pm0.004$ & $\mathbf{0.449}\pm0.003$ \\
    \ourmodel{} (pred.\ label)$^\dagger$
      & 0.821 & 0.825 & 1.438
      & 0.518 & 1.067 & 0.451 \\
    \midrule
    \multicolumn{7}{l}{\small\textit{Centre-bias analysis (seed~42)}} \\[2pt]
    Centre-bias upper bound$^\S$
      & 0.831 & 0.835 & 1.557 & 0.556 & 0.999 & 0.464 \\
    VideoSwin (debiased$^\ddagger$)
      & 0.669 & 0.671 & 0.726 & 0.270 & -- & -- \\
    \ourmodel{} (ours, debiased$^\ddagger$)
      & 0.765 & 0.769 & 1.188 & 0.431 & -- & -- \\
    \bottomrule
  \end{tabular}}
  \label{tab:sota}
\end{table*}

\ourmodel{} surpasses ACL-Net across all reported metrics:
$\Delta$AUC-J\,=\,+0.003, $\Delta$NSS\,=\,+0.184 (+14.7\%),
$\Delta$CC\,=\,+0.067 (+14.9\%).  Several aspects of this comparison
warrant detailed analysis.

\mypar{Asymmetric NSS vs.\ AUC-J gains.}
The most informative comparison is the relative magnitude of NSS and AUC-J
improvements.  FiLM conditioning over the VideoSwin baseline improves NSS
by $+9.2\%$ but AUC-J by only $+1.7\%$.  This asymmetry reflects the
fundamental difference between these metrics: AUC-J is a ranking metric
assessing whether fixation pixels receive higher predicted saliency than
non-fixation pixels---the VideoSwin backbone's temporally-aware features
already rank these pixels reasonably well.  NSS, by contrast, measures
the normalised predicted saliency value \emph{precisely at} the fixation
locations, quantifying how sharply and accurately the model fires at the
exact spatial positions observers fixated.  The large NSS gain with modest
AUC-J gain indicates that FiLM conditioning primarily improves the
\emph{calibrated precision} of saliency responses at fixation locations,
not the gross pixel ranking.  Mechanistically, this is consistent with the
FiLM modulation selectively amplifying channels encoding the type of visual
feature that is actually fixated in each density category---individual
silhouettes in SP, flow structures in DF, scene landmarks in DC---producing
sharper, more precisely placed peaks rather than simply reordering which
pixels score highest.  The $-5.8\%$ KL reduction (1.133~$\to$~1.068) and
$+8.6\%$ CC improvement corroborate this interpretation: both metrics are
sensitive to the shape and spatial structure of the saliency distribution,
not merely its ranking.

\mypar{VideoSwin backbone advantage.}
Comparing our VideoSwin baseline (AUC-J\,0.806, NSS\,1.313) against
ACL-Net (0.817, 1.250) reveals an interesting crossover: ACL-Net achieves
higher AUC-J ranking than our unconditioned baseline, suggesting that its
decoder-level density conditioning provides a useful spatial reranking
signal.  However, its NSS (1.250) is lower than our VideoSwin baseline
(1.313), meaning that despite better ranking, ACL-Net's predicted saliency
values are less precisely calibrated at actual fixation locations.
This reflects the representational advantage of Kinetics-400-pretrained
Swin3D-S for temporal motion encoding: the stronger backbone captures
temporally-resolved motion patterns that directly predict where observers
fixate, rather than only which areas are more salient than others in a
coarse sense.  The addition of FiLM conditioning then provides the precision
gains on top of this strong backbone foundation.

\mypar{Debiased evaluation reveals the true conditioning gain.}
Standard saliency metrics conflate crowd-content-responsive predictions with
centre-prior alignment: if a model learns to predict a fixed central blob
(matching the dataset's overall fixation tendency), it will score well on
NSS and AUC-J regardless of whether it understands crowd structure.
To isolate genuine crowd-specific contributions, we subtract the mean
training-set fixation density from all model predictions before computing
metrics (Table~\ref{tab:sota}, lower panel; Section~\ref{sec:sota}).

The centre-bias upper bound---predicting the mean training fixation density
for every test sample, with no model---achieves NSS\,1.557, exceeding both
VideoSwin (1.313) and \ourmodel{} (1.437) in standard evaluation.
This establishes that a substantial portion of all learned models' NSS
reflects centre-prior alignment rather than crowd-content understanding.

Under debiased evaluation the picture inverts dramatically.
VideoSwin's NSS collapses from 1.313 to 0.726, a 45\% loss: the
unconditioned Kinetics-pretrained backbone is heavily centre-biased, and
once the shared prior is subtracted, its crowd-specific predictions are weak.
\ourmodel{} drops from 1.437 to 1.188, a 17\% loss, because density
conditioning steers the bottleneck toward crowd-category-specific features
that are largely independent of the centre prior.
The result is a \emph{debiased conditioning gain} of
$\Delta\mathrm{NSS}\!=\!+0.462$ ($1.188 - 0.726$), which is
$3.7\!\times$ the standard gain of $+0.124$.

This finding reframes the contribution: density conditioning at the
bottleneck does not merely add precision on top of a centre-biased
backbone---it replaces the backbone's reliance on the centre prior with
predictions that are responsive to crowd density structure.

\subsection{Ablation Study}
\label{sec:ablation}

Table~\ref{tab:ablation} evaluates the contribution of each component.
All variants share the same VideoSwin backbone, 5-stage 2D decoder, and
hyperparameters; they differ only in what is added beyond the baseline.

\begin{table}[!t]
  \def\arraystretch{1.1}
  \centering
  \footnotesize
  \caption{%
    Ablation study on the CrowdFix test set (single seed, seed~42).
    $\Delta$NSS: gain over VideoSwin baseline (standard evaluation).
    $\Delta$NSS$_\text{deb}$: gain over debiased VideoSwin (NSS\,0.726);
    ``--'' where debiased eval was not run.
    The ``+Density FiLM'' row reports seed~42 values; Table~\ref{tab:sota}
    reports mean\,$\pm$\,std over 4 seeds.  \textbf{Bold}: best standard NSS.
    The bottom panel (Soft FiLM; Multiscale FiLM) reports Phase~2 variants
    that explore conditioning without GT labels at inference and multi-scale
    injection, respectively.
  }
  \resizebox{\columnwidth}{!}{%
  \begin{tabular}{lccccc}
    \toprule
    \textsc{Variant}
      & \textsc{AUC-J}
      & \textsc{NSS}
      & \textsc{CC}
      & $\Delta$\textsc{NSS}
      & $\Delta$\textsc{NSS}$_\text{deb}$ \\
    \midrule
    VideoSwin (baseline)
      & 0.806 & 1.313 & 0.476 & -- & -- \\
    +~Density FiLM (\textbf{ours})
      & \textbf{0.820} & \textbf{1.437} & \textbf{0.517} & $+$0.124 & $+$0.462 \\
    \midrule
    +~One-hot density code (no FiLM)
      & 0.819 & 1.424 & 0.514 & $+$0.111 & $+$0.345 \\
    +~Optical flow (RAFT)
      & 0.820 & 1.437 & 0.517 & $+$0.124 & -- \\
    +~3D decoder + social prior
      & 0.808 & 1.313 & 0.477 & $\pm$0.000 & -- \\
    \midrule
    +~Soft FiLM (self-sup.\ at inference)
      & 0.820 & 1.424 & 0.514 & $+$0.111 & $+$0.446 \\
    +~Multiscale FiLM (strides 32/16/8/4)
      & 0.819 & 1.412 & 0.512 & $+$0.099 & $+$0.405 \\
    +~Continuous (VGG16 features, no GT label)
      & \textbf{0.821} & \textbf{1.441} & \textbf{0.520} & $+$0.128 & $+$0.416 \\
    \bottomrule
  \end{tabular}}
  \label{tab:ablation}
\end{table}

\mypar{Density FiLM is the decisive and sufficient improvement.}
FiLM bottleneck conditioning improves NSS by $+0.124$ and CC by $+0.041$
over the plain VideoSwin baseline, with no other architectural changes.
The entire performance gain over the unconditioned backbone---and over
ACL-Net---is attributable to a single $64\!\times\!1536$ linear projection
and a $3\!\times\!64$ embedding matrix.  This confirms that crowd density
labels carry information about fixation patterns that is not implicit in
the video features, and that the bottleneck is the correct injection point
for this information to propagate through the entire decoding process.

\mypar{One-hot conditioning ablation.}
To isolate the contribution of the FiLM mechanism itself---as distinct from
simply having any density signal---we train a variant that injects a 3-class
one-hot vector via a direct $3\!\times\!768$ linear bias added to the
bottleneck, with no learned embedding and no multiplicative $\gamma$ pathway.
This variant, \textit{+One-hot density code}, uses the same backbone,
decoder, auxiliary classification head, and hyperparameters as \ourmodel{}
but replaces the $64$-dim embedding and FiLM projection with the 3-dim
one-hot bias (2,304 parameters vs.\ 100K for FiLM).

The one-hot model achieves NSS\,1.424, AUC-J\,0.819, CC\,0.514
(Table~\ref{tab:ablation}), versus \ourmodel's NSS\,1.437, AUC-J\,0.820,
CC\,0.517 (seed~42).  Two findings follow from standard evaluation.
First, the conditioning \emph{signal} is the primary driver: the one-hot model
captures $\Delta$NSS\,=\,$+0.111$ of FiLM's $+0.124$ total gain (89\%).
Second, the FiLM \emph{mechanism} adds $+0.013$ NSS
(${\approx}1.9\times$ the per-seed standard deviation of $0.007$),
a difference that is indicative but not definitive at conventional
significance thresholds; the near-orthogonal $\gamma$ geometry
(Section~\ref{sec:film_analysis}) provides geometrically independent support.

Debiased evaluation substantially sharpens this picture
(Table~\ref{tab:ablation}, $\Delta$NSS$_\text{deb}$ column).
The one-hot model's debiased NSS is 1.071 ($\Delta$NSS$_\text{deb}$\,=\,$+0.345$
over debiased VideoSwin 0.726), while \ourmodel{} reaches 1.188
($\Delta$NSS$_\text{deb}$\,=\,$+0.462$).  The FiLM advantage over one-hot
grows from $+0.013$ standard to $+0.117$ debiased---a $9\times$ amplification.
An additive one-hot bias shifts all bottleneck channels uniformly in a
class-specific direction; FiLM's multiplicative $\gamma$ pathway independently
scales each channel, allowing it to amplify channels that encode crowd-specific
features while suppressing channels that encode the centre prior.
This asymmetric channel-level control is precisely what survives the removal
of the shared centre-prior, explaining why the FiLM mechanism advantage is
far more visible after debiasing.

\mypar{Soft FiLM: self-supervised conditioning closes the gap.}
Rather than conditioning on GT density labels, the Soft FiLM variant derives
its conditioning signal from the model's own density head via
$\tilde{\mathbf{e}} = \text{softmax}(\hat{\mathbf{z}}_\text{det}) \mathbf{W}_\text{emb}$,
where $\hat{\mathbf{z}}_\text{det}$ are stop-gradient logits and
$\mathbf{W}_\text{emb}$ is the embedding weight matrix.
This enables continuous interpolation between the three density embeddings
and removes the GT label requirement at inference entirely.
Soft FiLM achieves NSS\,1.424, AUC-J\,0.820, CC\,0.514
($\Delta$NSS$_\text{deb}$\,=\,$+0.446$)---only $-0.016$ NSS below the oracle-label
\ourmodel{} ($+0.462$ debiased).  The near-identical debiased gain
demonstrates that the conditioning benefit comes primarily from the density
\emph{representation} in the embedding space, not from oracle-label precision.
Since Soft FiLM is fully self-supervised at inference, it is the preferred
deployment variant when GT crowd density labels are unavailable.

\mypar{Multiscale FiLM: bottleneck injection is the optimal point.}
Extending density conditioning to all four decoder stages (strides 32, 16, 8, 4)
with independent per-scale FiLM projections (zero-initialised to preserve
training stability) yields NSS\,1.412, AUC-J\,0.819
($\Delta$NSS$_\text{deb}$\,=\,$+0.405$)---weaker than both single-bottleneck
variants.  This result confirms that the bottleneck is already the optimal
conditioning injection point.  By the time features enter the decoder,
density-specific routing has already been established by the bottleneck FiLM;
applying further modulation at finer scales adds degrees of freedom that the
available training data cannot support, consistent with the variance--bias
analysis for the three-branch ablation below.

\mypar{Continuous VGG16 conditioning: highest oracle NSS, but learned embedding separates density better.}
Replacing the discrete 3-class embedding with a continuous density descriptor derived
from the frozen VGG16 frontend of CSRNet~\citep{csrnet} (conv1--conv4\_3, ImageNet-pretrained)
followed by dilated convolutional backend layers yields NSS\,1.441, AUC-J\,0.821,
CC\,0.520 ($\Delta$NSS$_\text{deb}$\,=\,$+0.416$) --- the highest \emph{standard} NSS
of all variants, but below the discrete FiLM model's debiased gain ($+0.462$).
Two findings emerge.  First, the continuous model achieves the best standard NSS
because VGG16 features encode scene-level properties correlated with the shared
centre-prior (blurriness, texture density, clutter), which boost fixation overlap
but do not isolate crowd-type-specific saliency.
Second, the debiased comparison reveals that the learned discrete embedding achieves
\emph{cleaner density-discriminative separation}: the discrete embedding is explicitly
supervised to distinguish the three density regimes, so its (γ, β) vectors encode
regime-specific channel recalibration that survives centre-prior removal.
The continuous VGG16 descriptor, by contrast, encodes a richer but less crowd-type-specific
signal that correlates more with centre-prior-compatible scene statistics.
This dissociation --- higher standard NSS, lower debiased NSS --- demonstrates that
VGG16 scene features and crowd-type-specific conditioning are complementary rather than
redundant, and motivates future work combining both in a dual-stream conditioning pathway.

\mypar{RAFT optical flow is redundant.}
Appending RAFT~\citep{raft} dense optical flow as a second input stream,
fused with the Swin3D bottleneck via cross-attention, yields metrics
identical to \ourmodel{} across all three metrics (AUC-J\,0.820,
NSS\,1.437, CC\,0.517; both seed~42).  This null result is informative rather than
disappointing: it demonstrates that the shifted 3D window attention in
Swin3D-S already captures the temporal dynamics that are relevant for
saliency prediction.  The Swin3D encoder computes attended temporal
correspondences across all four hierarchical stages at multiple spatial
scales; by the final bottleneck, every spatial token has been shaped by
token-level temporal attention patterns encoding implicit velocity and
trajectory information.  RAFT optical flow, which computes instantaneous
per-pixel displacement vectors from adjacent frame pairs, encodes a
structurally redundant and representationally shallower form of the same
motion signal.  The cross-attention fusion layer learns to zero-gate the
RAFT stream, contributing nothing beyond what the backbone already knows,
while adding $\approx$10M parameters that impose a slight regularisation
burden.

\mypar{Added capacity cancels FiLM and collapses to baseline.}
The three-branch architecture---a parallel UpBlock3d temporal decoder
branch, a per-category spatial social force prior, and a learned
gate fusing the FiLM and 3D branches---achieves AUC-J\,0.808,
NSS\,1.313, CC\,0.477.  The three-branch NSS of $\approx$1.313 is the most striking result in
the ablation: it is effectively identical to the VideoSwin baseline
($\Delta$NSS\,$\approx$\,0), meaning the additional architecture not only
fails to improve over the FiLM model but also cancels the FiLM gain
entirely.  This outcome is a direct consequence of the variance--bias
tradeoff at the available training scale.  The three-branch model
introduces sufficient estimation variance---from the 3D decoder's
additional parameters, the per-category social force spatial bases, and
the fusion gate---that the expected reduction in approximation bias from
these priors is negated.  In effect, the model cannot determine which
branch to trust on 3,189 training clips; both branches push toward zero
activation and the FiLM modulation is neutralised.  This result provides
a qualitative upper bound on the model complexity that CrowdFix's
training scale can support: architecturally, any addition beyond the
FiLM layer falls outside the data-supportable complexity budget.

\subsection{Per-Density-Category Performance}
\label{sec:per_cat}

Table~\ref{tab:per_cat} breaks down \ourmodel{} performance by density
category, providing insight into how density conditioning benefits each
regime.

\begin{table}[!t]
  \def\arraystretch{1.1}
  \centering
  \footnotesize
  \caption{%
    Per-density-category test performance of \ourmodel{} (seed~42).
    Overall row corresponds to the seed~42 run; Table~\ref{tab:sota}
    reports multi-seed mean\,$\pm$\,std.
    $n$\,=\,number of test clips.
  }
  \begin{tabular}{lccc}
    \toprule
    \textsc{Category} ($n$)
      & \textsc{AUC-J}$\uparrow$
      & \textsc{NSS}$\uparrow$
      & \textsc{CC}$\uparrow$ \\
    \midrule
    SP -- Sparse Pedestrian (1443)  & 0.821 & 1.439 & 0.523 \\
    DF -- Dense Free-flowing (1924) & 0.818 & 1.422 & 0.514 \\
    DC -- Dense Congested (1989)    & 0.822 & 1.449 & 0.517 \\
    \midrule
    \textit{Overall} (5356)         & 0.820 & 1.437 & 0.517 \\
    \bottomrule
  \end{tabular}
  \label{tab:per_cat}
\end{table}

DC achieves the highest NSS (1.449) despite representing the most
challenging density scenario.  We interpret this as the strongest evidence
for the effectiveness of density conditioning: in congested scenes, the
density label is uniquely informative.  Without the DC conditioning signal,
the model would be uncertain whether to deploy an individual-tracking or
a collective-structure strategy; the DC label unambiguously signals that
neither is appropriate and that scene-level landmarks should be prioritised.
The corresponding FiLM parameters ($\gamma_\mathrm{DC}$, $\beta_\mathrm{DC}$)
must have learned to amplify channels encoding high-contrast edge detection
and semantic region boundaries, which precisely corresponds to the fixation
targets in congested subway and corridor scenes (bright signage, architectural
apertures).

SP achieves the highest CC (0.523), reflecting that sparse crowd saliency
distributions are spatially simpler: they consist of a small number of
well-separated peaks centred on individual bodies, yielding a smooth,
approximately unimodal distribution that is relatively easy to fit with
high correlation.  The SP conditioning directs the model toward
spot-attention behaviour---amplifying channels selective for localised
moving-object boundaries---producing sharp, well-placed saliency peaks.

The SP category spans a wide range of group sizes, from single pedestrians
to groups of 10--20 persons in open plazas.  Fixation strategy is expected
to vary within SP accordingly---individual trajectory tracking for lone
figures, social gaze-following for small groups, and broad scene scanning
for larger gatherings---yet the per-category NSS above does not resolve
this variation.  Disaggregating SP performance by approximate group size
using existing test-set predictions is a planned zero-cost analysis.

DF achieves the lowest scores across all metrics, consistent with the
inherent difficulty of the free-flowing regime: ground-truth fixations are
broadly distributed over multiple equivalently-attractive individuals, yielding
a diffuse GT saliency map with no dominant peak.  Predicting diffuse
distributions with high NSS is intrinsically harder than predicting peaked
ones, since the normalised saliency value at any individual fixation location
is lower when fixations are spread across many locations.  This suggests
that DF may require finer temporal modelling of specific flow-trajectory
emergence points to achieve the per-fixation precision that NSS rewards.

\subsection{FiLM Parameter Analysis}
\label{sec:film_analysis}

A potential concern is that with only $K\!=\!3$ density classes, the
$3\!\times\!64$ embedding table degenerates to three fixed $(\gamma,\beta)$
vectors---a class-conditional bias lookup rather than dynamic feature
modulation.  We verify whether the learned parameters encode genuinely
distinct structure by extracting the three $(\gamma,\beta)\!\in\!\mathbb{R}^{768}$
vectors from the trained checkpoint and computing pairwise cosine similarity.

The $\gamma$ vectors are nearly orthogonal: SP\,vs.\,DF $= -0.03$,
SP\,vs.\,DC $= -0.10$, DF\,vs.\,DC $= -0.06$.  Critically, negative
cosine similarity means the SP scale factors are anti-correlated with
those of DC: channels amplified in sparse scenes are suppressed in
congested scenes and vice versa.  This is precisely the structure
that would be expected if the model has learned to spotlight
individual-body channels for SP and scene-boundary channels for DC.
The maximum pairwise similarity of $-0.03 \ll 0.99$ confirms that
the embedding is non-degenerate.

A channel-level inspection confirms the asymmetry: the single most
discriminative channel (Ch\,437) exhibits $\gamma_\mathrm{SP}\!=\!+1.09$
vs.\ $\gamma_\mathrm{DC}\!=\!-1.88$ (range\,2.97), while channel 291
shows the complementary pattern ($\gamma_\mathrm{SP}\!=\!-2.54$,
$\gamma_\mathrm{DC}\!=\!+0.14$).  The L2 norms of the full $(\gamma,\beta)$
vectors are comparable across classes (SP\,25.3, DF\,25.9, DC\,24.9),
confirming that all three density types receive equally strong modulation
of the same magnitude---balanced but geometrically distinct.

\subsection{Qualitative Analysis}

Figure~\ref{fig:qualitative} presents a stratified sample of model
behaviour across each density category: one best-case frame (top NSS
in the test split), one average-case frame (NSS near the per-category
median of $\approx$1.43), and one challenging frame (negative NSS,
indicating saliency placed away from observer fixations).  All nine
frames are from distinct videos.  Each image combines the original
frame, the predicted saliency overlay (INFERNO colormap with
saliency-proportional alpha), and per-observer fixation rings (white
ring with black halo at 2-degree visual-angle radius; white centre dot).

\mypar{Best-case frames (left column, NSS $\approx$ 3.8--4.3).}
In all three categories the model correctly matches observer fixation
strategy: spot attention on proximal individuals (SP video~212,
frame~43, NSS\,=\,3.83); diffuse coverage over the pedestrian flow
(DF video~152, frame~36, NSS\,=\,3.92); and precise co-localisation
with a high-contrast scene landmark rather than crowd mass
(DC video~335, frame~50, NSS\,=\,4.34).  These frames share strong
centre alignment between fixation clusters and predicted peaks,
partly attributable to dataset centre bias but primarily driven by
the density conditioning placing saliency at structurally salient
regions regardless of image position.

\mypar{Average-case frames (middle column, NSS $\approx$ 1.4).}
Performance near the per-category mean reveals a broader, less
precise saliency distribution.  Predicted peaks partially overlap
the fixation cluster but with more spatial spread, consistent with
the model averaging over plausible fixation locations rather than
committing to the correct target.  The fixation rings are still
largely enclosed by the saliency blob, confirming above-chance
localisation at the cost of precision.

\mypar{Challenging frames (right column, NSS $<$ 0).}
Negative NSS indicates the model actively placed probability mass
away from where observers looked.  In all three cases the root cause
is centre bias: the model predicts a central saliency peak while the
observer fixations are displaced to the periphery --- an off-centre
individual (SP video~260, frame~81), a lateral flow structure
(DF video~093, frame~17), or an off-centre landmark
(DC video~180, frame~38).  These failure cases confirm that centre
bias inherited from the pre-trained VideoSwin encoder is the primary
remaining limitation of the model.

\begin{figure*}[p]
  \centering
  \setlength{\tabcolsep}{3pt}
  \begin{tabular}{c@{\hspace{4pt}}ccc}
    & \textbf{Best case} & \textbf{Average case} & \textbf{Challenging case} \\[4pt]
    \rotatebox{90}{\parbox{2.8cm}{\centering\textbf{SP}\\\small\textit{Sparse Ped.}}} &
      \includegraphics[width=0.305\textwidth]{../results/212/frame_00043.png} &
      \includegraphics[width=0.305\textwidth]{../results/163/frame_00042.png} &
      \includegraphics[width=0.305\textwidth]{../results/260/frame_00081.png} \\
    \multicolumn{1}{c}{} &
      \small v212\,f43 \ NSS\,=\,3.83 &
      \small v163\,f42 \ NSS\,=\,1.44 &
      \small v260\,f81 \ NSS\,=\,$-$0.20 \\[6pt]
    \rotatebox{90}{\parbox{2.8cm}{\centering\textbf{DF}\\\small\textit{Dense Free-flow}}} &
      \includegraphics[width=0.305\textwidth]{../results/152/frame_00036.png} &
      \includegraphics[width=0.305\textwidth]{../results/064/frame_00080.png} &
      \includegraphics[width=0.305\textwidth]{../results/093/frame_00017.png} \\
    \multicolumn{1}{c}{} &
      \small v152\,f36 \ NSS\,=\,3.92 &
      \small v064\,f80 \ NSS\,=\,1.43 &
      \small v093\,f17 \ NSS\,=\,$-$0.43 \\[6pt]
    \rotatebox{90}{\parbox{2.8cm}{\centering\textbf{DC}\\\small\textit{Dense Congested}}} &
      \includegraphics[width=0.305\textwidth]{../results/335/frame_00050.png} &
      \includegraphics[width=0.305\textwidth]{../results/160/frame_00014.png} &
      \includegraphics[width=0.305\textwidth]{../results/180/frame_00038.png} \\
    \multicolumn{1}{c}{} &
      \small v335\,f50 \ NSS\,=\,4.34 &
      \small v160\,f14 \ NSS\,=\,1.47 &
      \small v180\,f38 \ NSS\,=\,$-$0.40 \\
  \end{tabular}
  \caption{%
    Qualitative results stratified by performance tier (nine unique videos).
    Each image overlays the predicted saliency (INFERNO colormap,
    saliency-proportional alpha) and per-observer fixation rings
    (2-degree visual-angle radius; white ring with black halo;
    white centre dot).
    \textbf{Best-case (left):} predicted saliency peaks co-localise
    tightly with observer fixations --- spot-attention on individuals (SP),
    diffuse flow coverage (DF), and landmark-level precision (DC).
    \textbf{Average-case (centre):} predictions partially overlap fixation
    clusters with broader spatial spread, consistent with
    above-chance localisation at median performance.
    \textbf{Challenging (right):} negative NSS indicates misplaced
    saliency --- the model predicts a central peak while fixations are
    off-centre (an eccentric individual in SP, a lateral flow structure
    in DF, an off-centre landmark in DC), exposing residual centre bias
    from the pre-trained encoder.
  }
  \label{fig:qualitative}
\end{figure*}

% --------------------------------------------------------------------------
\section{Conclusion}
\label{sec:conclusion}
% --------------------------------------------------------------------------

We presented \ourmodel, a Video Swin Transformer architecture augmented with
a FiLM-conditioned bottleneck for crowd density-aware video saliency
prediction.  The central claim of this work---that crowd density induces
categorical fixation regime shifts that require explicit conditioning to
model accurately---is supported by a consistent body of evidence: the FiLM
module's $\approx$100K parameters improve NSS by $+9.2\%$ and CC by $+8.6\%$
over the unconditioned backbone in standard evaluation, while three
progressively more complex extensions all fail to improve, with the most
complex variant collapsing NSS precisely to baseline.  The result is a new
state of the art on CrowdFix
(AUC-J\,$0.820\pm0.001$, NSS\,$1.434\pm0.007$, CC\,$0.517\pm0.003$, mean over 4~seeds),
representing the first published improvement over ACL-Net (2019) in six years.

Debiased evaluation strengthens this conclusion.  Subtracting the mean
training fixation density from all predictions to remove shared centre-prior
alignment reveals that the unconditioned VideoSwin backbone loses 45\% of
its NSS ($1.313\to0.726$) while \ourmodel{} loses only 17\%
($1.437\to1.188$), yielding a debiased conditioning gain of
$\Delta\mathrm{NSS}\!=\!+0.462$---$3.7\!\times$ the standard gain.
Density conditioning does not merely sharpen an already centre-biased
response; it replaces centre-prior exploitation with crowd-category-specific
feature activation that survives the removal of the shared prior.

The quantitative characterisation of the capacity--regularisation boundary
on CrowdFix has broader implications.  The precise NSS collapse
($\Delta\mathrm{NSS}\!=\!0.000$ for the three-branch model) establishes
3,189 training clips as a threshold below which only minimal targeted
inductive bias is data-supportable, providing a concrete empirical reference
for future work on small-scale crowd video benchmarks.  Per-category
analysis reveals that DC scenes benefit most from density conditioning
(NSS\,1.449), confirming that the label is most informative precisely in
the regime where crowd structure is most ambiguous to a
density-agnostic model.

Several directions remain open.  Self-supervised pretraining on large
unlabelled crowd video corpora would expand the effective training scale and
potentially allow the complex temporal architectures shown to overfit here
to achieve their intended representational benefit.  Continuous density
conditioning---replacing the 3-class label with a learned embedding of
per-frame local density estimates---could capture within-category variation
that the coarse label discards.  Finally, extending the dataset with
additional density categories (e.g., extreme evacuation dynamics) and
longer eye-tracking sequences per clip would provide the data density
required for finer-grained conditioning schemes.

\subsection*{Limitations}

Three limitations bound the scope of this work.
First, \textbf{oracle density labels}: at inference the model can use
either the ground-truth CrowdFix category label or its own auxiliary
density head prediction.  Empirically, the predicted-label variant
achieves NSS\,1.438 vs.\ oracle NSS\,1.437 on seed~42 ($\Delta\!=\!0.001$;
see Table~\ref{tab:sota}), confirming that the density head classifies
crowd categories accurately enough that oracle access is not required in
practice.  The remaining limitation is that the SP/DF/DC taxonomy is
CrowdFix-specific; generalising to an unlabelled corpus would require
either a compatible density classifier or unsupervised cluster assignment.
Second, \textbf{single-dataset scope}: all results are on CrowdFix.
Whether the density-conditioned bottleneck generalises to other crowd
video collections (e.g.\ crowd scenes in DHF1K or Hollywood-2) has not
been tested.
Third, \textbf{centre bias and debiased evaluation}: the VideoSwin encoder
inherits a centre-biased prior from Kinetics-400 pretraining, and the
CrowdFix fixation maps exhibit a corresponding centre tendency
(negative-NSS examples in Figure~\ref{fig:qualitative} are all off-centre
misses).  Standard metrics therefore conflate crowd-content-responsive
predictions with centre-prior alignment.  Debiased evaluation
(Section~\ref{sec:sota}) reveals that this conflation favours the
unconditioned backbone more than it favours \ourmodel{}: VideoSwin loses
45\% of its NSS under debiasing (1.313~$\to$~0.726) while \ourmodel{} loses
only 17\% (1.437~$\to$~1.188).  The centre-bias upper bound (NSS\,1.557)
sets a ceiling no learned model has yet exceeded, representing an open
challenge for future work on bias-corrected crowd video saliency.

\section*{Acknowledgements}

The author acknowledges CSC -- IT Center for Science, Finland, for
computational resources on the LUMI HPC cluster under project~462000131.

% --- Bibliography ---
\bibliographystyle{unsrtnat}
\bibliography{references}

\end{document}
